% ============================================================================
% DRAFTING HEADER (strip before deposit)
% R904 (Fable 5.1, Nov 06 2026), at the author's direction ("Would you please now write a paper on psychology for me
% using this note of mine here from today"): the author's note of November 6, 2026, 10:30 AM, converted into this
% paper. The note is quoted whole in Appendix A, with one clause corrected for grammar at his direction the same day
% (the original clause is kept in a comment there). The case, the root claim, the double bind, the Cassandra and the
% voice formulations are the author's; the formal model, the placement in the literature and the predictions are the
% conversion's, written to carry his claim. The seventh extraction of the corpus and the second psychological one,
% beside norm_cost.tex. References checked against the sources on Nov 6, 2026, except where a comment says otherwise.
% No real name of any person in the case is used; the note's own redaction is kept.
% R906 (Fable 5.1, Nov 06 2026): revised on the first referee's brief (ChatGPT, Nov 6 2026); the R905 body is kept in r906/unconveyable_pre_r906.tex.
% ============================================================================
\begin{filecontents}[overwrite]{unconveyable_refs.bib}
@misc{Aethus, author={Benander, Ajax}, title={Aethic Reasoning: A Comprehensive Solution to the Quantum Measurement Problem}, year={2024}, month = {11}, url={https://philpapers.org/rec/BENARA-5}, note={Manuscript; the complete development of the framework, successive deposits November 2024--2026. This version: November 2026.}}
@misc{NormCost,
  author = {Benander, Ajax},
  title = {Cognitive Traits Across Domains: Classification, Symmetry and the Cost of Interaction},
  year = {2026}, month = {10},
  howpublished = {Working draft},
  note = {The companion psychological extraction of the same framework: the Branch Principle, the Principle of Symmetry, and a cost model of interaction under governing norms.}}
@article{Milton2012,
  author = {Milton, Damian E. M.},
  title = {On the ontological status of autism: the `double empathy problem'},
  journal = {Disability \& Society}, volume = {27}, number = {6}, pages = {883--887}, year = {2012},
  doi = {10.1080/09687599.2012.710008}}
@article{Chevallier2012,
  author = {Chevallier, Coralie and Kohls, Gregor and Troiani, Vanessa and Brodkin, Edward S. and Schultz, Robert T.},
  title = {The social motivation theory of autism},
  journal = {Trends in Cognitive Sciences}, volume = {16}, number = {4}, pages = {231--239}, year = {2012},
  doi = {10.1016/j.tics.2012.02.007}}
@article{Jaswal2019,
  author = {Jaswal, Vikram K. and Akhtar, Nameera},
  title = {Being versus appearing socially uninterested: Challenging assumptions about social motivation in autism},
  journal = {Behavioral and Brain Sciences}, volume = {42}, pages = {e82}, year = {2019},
  doi = {10.1017/S0140525X18001826}}
@article{Ridgway2025,
  author = {Ridgway, Kathryn and Cooke, Kassandrah and Demmer, David H. and Hooley, Merrilyn and Westrupp, Elizabeth and Stokes, Mark A.},
  title = {Camouflaging autism in pursuit of friendship and intimate relationships: A systematic review},
  journal = {Autism in Adulthood}, volume = {7}, number = {5}, pages = {543--557}, year = {2025},
  doi = {10.1089/aut.2023.0160}}
@article{Hull2017,
  author = {Hull, Laura and Petrides, K. V. and Allison, Carrie and Smith, Paula and Baron-Cohen, Simon and Lai, Meng-Chuan and Mandy, William},
  title = {`Putting on my best normal': Social camouflaging in adults with autism spectrum conditions},
  journal = {Journal of Autism and Developmental Disorders}, volume = {47}, number = {8}, pages = {2519--2534}, year = {2017},
  doi = {10.1007/s10803-017-3166-5}}
@article{Hull2019,
  author = {Hull, Laura and Mandy, William and Lai, Meng-Chuan and Baron-Cohen, Simon and Allison, Carrie and Smith, Paula and Petrides, K. V.},
  title = {Development and validation of the {Camouflaging Autistic Traits Questionnaire} ({CAT-Q})},
  journal = {Journal of Autism and Developmental Disorders}, volume = {49}, number = {3}, pages = {819--833}, year = {2019},
  doi = {10.1007/s10803-018-3792-6}}
@article{Cage2019,
  author = {Cage, Eilidh and Troxell-Whitman, Zoe},
  title = {Understanding the reasons, contexts and costs of camouflaging for autistic adults},
  journal = {Journal of Autism and Developmental Disorders}, volume = {49}, number = {5}, pages = {1899--1911}, year = {2019},
  doi = {10.1007/s10803-018-03878-x}}
@article{Livingston2019,
  author = {Livingston, Lucy Anne and Shah, Punit and Happ{\'e}, Francesca},
  title = {Compensatory strategies below the behavioural surface in autism: a qualitative study},
  journal = {The Lancet Psychiatry}, volume = {6}, number = {9}, pages = {766--777}, year = {2019},
  doi = {10.1016/S2215-0366(19)30224-X}}
@article{Pearson2021,
  author = {Pearson, Amy and Rose, Kieran},
  title = {A conceptual analysis of autistic masking: Understanding the narrative of stigma and the illusion of choice},
  journal = {Autism in Adulthood}, volume = {3}, number = {1}, pages = {52--60}, year = {2021},
  doi = {10.1089/aut.2020.0043}}
@article{Botha2020,
  author = {Botha, Monique and Frost, David M.},
  title = {Extending the minority stress model to understand mental health problems experienced by the autistic population},
  journal = {Society and Mental Health}, volume = {10}, number = {1}, pages = {20--34}, year = {2020},
  doi = {10.1177/2156869318804297}}
@article{Raymaker2020,
  author = {Raymaker, Dora M. and Teo, Alan R. and Steckler, Nicole A. and Lentz, Brandy and Scharer, Mirah and Delos Santos, Austin and Kapp, Steven K. and Hunter, Morrigan and Joyce, Andee and Nicolaidis, Christina},
  title = {\mkbibquote{Having all of your internal resources exhausted beyond measure and being left with no clean-up crew}: Defining autistic burnout},
  journal = {Autism in Adulthood}, volume = {2}, number = {2}, pages = {132--143}, year = {2020},
  doi = {10.1089/aut.2019.0079}}
@article{Higgins2021,
  author = {Higgins, Julianne M. and Arnold, Samuel R. C. and Weise, Janelle and Pellicano, Elizabeth and Trollor, Julian N.},
  title = {Defining autistic burnout through experts by lived experience: Grounded {Delphi} method investigating \#{AutisticBurnout}},
  journal = {Autism}, volume = {25}, number = {8}, pages = {2356--2369}, year = {2021},
  doi = {10.1177/13623613211019858}}
@article{Mantzalas2022a,
  author = {Mantzalas, Jane and Richdale, Amanda L. and Dissanayake, Cheryl},
  title = {A conceptual model of risk and protective factors for autistic burnout},
  journal = {Autism Research}, volume = {15}, number = {6}, pages = {976--987}, year = {2022},
  doi = {10.1002/aur.2722}}
% volume, issue and pages of Mantzalas2022a from memory; the DOI checked.
@article{Mantzalas2022b,
  author = {Mantzalas, Jane and Richdale, Amanda L. and Adikari, Achini and Lowe, Jennifer and Dissanayake, Cheryl},
  title = {What is autistic burnout? {A} thematic analysis of posts on two online platforms},
  journal = {Autism in Adulthood}, volume = {4}, number = {1}, pages = {52--65}, year = {2022},
  doi = {10.1089/aut.2021.0021}}
@article{Crompton2020,
  author = {Crompton, Catherine J. and Ropar, Danielle and Evans-Williams, Claire V. M. and Flynn, Emma G. and Fletcher-Watson, Sue},
  title = {Autistic peer-to-peer information transfer is highly effective},
  journal = {Autism}, volume = {24}, number = {7}, pages = {1704--1712}, year = {2020},
  doi = {10.1177/1362361320919286}}
@article{Crompton2020b,
  author = {Crompton, Catherine J. and Hallett, Sonny and Ropar, Danielle and Flynn, Emma and Fletcher-Watson, Sue},
  title = {\mkbibquote{I never realised everybody felt as happy as I do when I am around autistic people}: A thematic analysis of autistic adults' relationships with autistic and neurotypical friends and family},
  journal = {Autism}, volume = {24}, number = {6}, pages = {1438--1448}, year = {2020},
  doi = {10.1177/1362361320908976}}
@article{Foster2025,
  author = {Foster, Sarah J. and Ackerman, Robert A. and Wilks, Charlotte E. H. and Dodd, Michelle and Calderon, Rachel and Ropar, Danielle and Fletcher-Watson, Sue and Crompton, Catherine J. and Sasson, Noah J.},
  title = {Rapport in same and mixed neurotype groups of autistic and non-autistic adults},
  journal = {Autism}, volume = {29}, number = {7}, year = {2025},
  doi = {10.1177/13623613251320444}}
@article{Sasson2017,
  author = {Sasson, Noah J. and Faso, Daniel J. and Nugent, Jack and Lovell, Sarah and Kennedy, Daniel P. and Grossman, Ruth B.},
  title = {Neurotypical peers are less willing to interact with those with autism based on thin slice judgments},
  journal = {Scientific Reports}, volume = {7}, pages = {40700}, year = {2017},
  doi = {10.1038/srep40700}}
@article{Silver2025,
  author = {Silver, Kate and Parsons, Sarah},
  title = {Strategies for developing interactional expertise: how non-autistic communication partners enable or block the effective contributions to significant conversations by autistic people},
  journal = {Disability and Rehabilitation}, year = {2025},
  note = {Published online 4 July 2025},
  doi = {10.1080/09638288.2025.2527940}}
@article{Black2023,
  author = {Black, Melissa H. and Clarke, Patrick J. F. and Deane, Evangeline and Smith, Dayna and Wiltshire, Grace and Yates, Emily and Lawson, Wenn B. and Chen, Nigel T. M.},
  title = {\mkbibquote{That impending dread sort of feeling}: Experiences of social interaction from the perspectives of autistic adults},
  journal = {Research in Autism Spectrum Disorders}, volume = {101}, pages = {102090}, year = {2023},
  doi = {10.1016/j.rasd.2022.102090}}
@article{Catala2021,
  author = {Catala, Amandine and Faucher, Luc and Poirier, Pierre},
  title = {Autism, epistemic injustice, and epistemic disablement: a relational account of epistemic agency},
  journal = {Synthese}, volume = {199}, pages = {9013--9039}, year = {2021},
  doi = {10.1007/s11229-021-03192-7}}
@book{Fricker2007,
  author = {Fricker, Miranda},
  title = {Epistemic Injustice: Power and the Ethics of Knowing},
  publisher = {Oxford University Press}, address = {Oxford}, year = {2007}}
@article{Sacks1974,
  author = {Sacks, Harvey and Schegloff, Emanuel A. and Jefferson, Gail},
  title = {A simplest systematics for the organization of turn-taking for conversation},
  journal = {Language}, volume = {50}, number = {4}, pages = {696--735}, year = {1974}}
@article{Gilbert1995,
  author = {Gilbert, Daniel T. and Malone, Patrick S.},
  title = {The correspondence bias},
  journal = {Psychological Bulletin}, volume = {117}, number = {1}, pages = {21--38}, year = {1995}}
@article{Maier1976,
  author = {Maier, Steven F. and Seligman, Martin E. P.},
  title = {Learned helplessness: Theory and evidence},
  journal = {Journal of Experimental Psychology: General}, volume = {105}, number = {1}, pages = {3--46}, year = {1976}}
@article{MacDonald1993,
  author = {MacDonald, J. F.},
  title = {Russell, {Wittgenstein} and the problem of the rhinoceros},
  journal = {The Southern Journal of Philosophy}, volume = {31}, number = {4}, pages = {409--424}, year = {1993}}
@book{Aeschylus,
  author = {Aeschylus},
  title = {Oresteia: Agamemnon, The Libation Bearers, The Eumenides},
  translator = {Lattimore, Richmond},
  publisher = {University of Chicago Press}, address = {Chicago}, year = {1953},
  note = {In \emph{Aeschylus I}, The Complete Greek Tragedies; the Cassandra scene of the \emph{Agamemnon}}}
@article{SassonMorrison2019,
  author = {Sasson, Noah J. and Morrison, Kerrianne E.},
  title = {First impressions of adults with autism improve with diagnostic disclosure and increased autism knowledge of peers},
  journal = {Autism}, volume = {23}, number = {1}, pages = {50--59}, year = {2019},
  doi = {10.1177/1362361317729526}}
@article{Crompton2025,
  author = {Crompton, Catherine J. and Foster, Sarah J. and Wilks, Charlotte E. H. and Dodd, Michelle and Efthimiou, Themis N. and Ropar, Danielle and Sasson, Noah J. and Lages, Martin and Fletcher-Watson, Sue},
  title = {Information transfer within and between autistic and non-autistic people},
  journal = {Nature Human Behaviour}, volume = {9}, number = {7}, pages = {1488--1500}, year = {2025},
  doi = {10.1038/s41562-025-02163-z}}
@article{Dotson2011,
  author = {Dotson, Kristie},
  title = {Tracking epistemic violence, tracking practices of silencing},
  journal = {Hypatia}, volume = {26}, number = {2}, pages = {236--257}, year = {2011},
  doi = {10.1111/j.1527-2001.2011.01177.x}}
@article{Mantzalas2024,
  author = {Mantzalas, Jane and Richdale, Amanda L. and Li, Xia and Dissanayake, Cheryl},
  title = {Measuring and validating autistic burnout},
  journal = {Autism Research}, volume = {17}, number = {7}, pages = {1417--1449}, year = {2024},
  doi = {10.1002/aur.3129}}
% R906 (Fable 5.1, Nov 06 2026), at the author's direction: four references added on the first referee's brief (the disclosure study; the registered replication of the diffusion-chain study; Dotson on silencing; the burnout measurement study). Sasson and Morrison's volume and pages from the publisher's record; Crompton et al. 2025 from the journal's record; Dotson from the journal's record; Mantzalas et al. 2024 from the journal's record.
\end{filecontents}

\documentclass[11pt]{article}
\usepackage[margin=1.1in]{geometry}
\usepackage[utf8]{inputenc}
\usepackage{amsmath, amssymb, amsthm}
\usepackage[hidelinks]{hyperref}
\usepackage{microtype}
\usepackage{booktabs}
\usepackage{array}
\usepackage{enumitem}
\usepackage[style=numeric, sorting=none, backend=bibtex]{biblatex}
\addbibresource{unconveyable_refs.bib}
\newcommand{\propositionautorefname}{Proposition}
\newcommand{\definitionautorefname}{Definition}
\newcommand{\hypothesisautorefname}{Hypothesis}
\newcommand{\corollaryautorefname}{Corollary}
\newcommand{\remarkautorefname}{Remark}
\renewcommand{\sectionautorefname}{Section}
\renewcommand{\subsectionautorefname}{Section}

\newtheorem{proposition}{Proposition}
\newtheorem{corollary}{Corollary}
\theoremstyle{definition}
\newtheorem{definition}{Definition}
\newtheorem{hypothesis}{Hypothesis}
\newtheorem{assumption}{Premise}
\newtheorem{remark}{Remark}
\newcommand{\E}{\mathbb{E}}
\newcommand{\PP}{\mathbb{P}}

\title{The Unconveyable Situation:\\ \large Obstructed Conveyance, Second-Order Silence,\\ \large and a Single-Root Hypothesis for Autistic Masking}
\author{Ajax Benander\thanks{Email: \href{mailto:akb9408@rit.edu}{akb9408@rit.edu}, Website: \url{https://aethus.org}. The author wrote the notes on which this paper rests on November 6, 2026 (Appendix~\ref{app:note}), and had Claude Fable 5.1 convert them into the paper at his direction the same day; the division between what is his and what is the conversion's is stated there, and the revision after the first referee's brief is recorded in Appendix~\ref{app:assumptions}. The paper is self-contained; its one companion in the psychology of interaction is a cost model of interaction under governing norms \cite{NormCost}, drawn from the same larger framework \cite{Aethus}, whose metaphysics plays no role here.}}
\date{November 6, 2026 --- working draft, revised}

\begin{document}
\maketitle

\begin{abstract}
The paper proposes a root for autistic masking: a pervasive inability to convey the truth of one's situation to others, which closes on itself. It starts from a dated first-person case, an autistic adult interrupted repeatedly while telling his father a story and then sent, with the words ``I don't want to interrupt but,'' to help his sister with a task. The case is read as an instance of a structure with three parts: contributions are obstructed; the account of the obstruction receives less uptake than the contributions it explains; and the resulting silence is read as a trait of the person, which strengthens the discount on the account. The paper's immediate contribution is this recursive mechanism, stated as a hypothesis about how masking and withdrawal are maintained in rooms where the three conditions hold; the author's stronger claim, that every masking effect and masking's contribution to burnout is an output of the one structure, is kept as an explicit hypothesis with its own boundaries and tests. An illustrative model with a Poisson process of bids for the floor shows, under stated assumptions about costs, policies and audience response, that the optimal policy resists interruptions only near a contribution's end, that a contribution is attempted only below a finite exposure which falls with its length, that the account of an obstruction has no higher modeled success than what it explains when it is no shorter and credited no better, that withdrawal is self-sustaining in the closed loop and left only by a change of room, channel or frame from elsewhere, and that at the resulting fixed point a definite part of the person's situation is closed under description, so that its own edge cannot be told. Masking is derived as the only speech option with positive expected net value in a band of exposures, and its per-episode cost enters an explicitly unvalidated burden bookkeeping. The account is placed beside the double empathy problem, social motivation theory, the camouflaging and burnout literatures, and the epistemic injustice literature, with the findings that qualify simple mismatch and symmetry claims included; the reply such an account most often receives, that the listener too has things never said, is shown to conflate a silence nothing holds in place with one the room holds in place; and eight studies are specified with what would weaken each claim.
\end{abstract}

\section{Introduction}

\subsection{A case}

On the morning of November 6, 2026, the author was telling his father a story, the one in which Russell, in the autumn of 1911, could not get the young Wittgenstein to admit that it was certain there was no rhinoceros in the room \cite{MacDonald1993}. By his count, a visitor to the house interrupted him about six times in two minutes. At the end she touched him on the shoulder and said, ``I don't want to interrupt but why don't you go help your sister download an app she's looking for.'' His note of 10:30 AM, written within the hour and quoted whole in Appendix~\ref{app:note}, reads the episode as the most concise root explanation of autistic masking he has managed to write, and the present paper is that note carried out. What the note reports, what the paper infers from it, and what it hypothesizes are kept apart throughout, beginning in \autoref{sec:case}.

The reading has four parts. First, what the visitor did was not to forbid the story but to keep constant pressure on him to mask, so that conveying the thought required consciously resisting, interruption by interruption. Second, a room full of allistic people scales the same pressure by the number of people in it, at which point the thought is not resisted for but given up: one learns to stay quiet. Third, the giving-up has a shape. Thinking of a funny story comes to translate, by a route the person can watch operate, into the conclusion that one had better keep one's mouth shut, so the isolation is produced inside the person by the room's pressure and then felt as one's own. Fourth, the result is read from outside as a deficit. The person is criticized for not making friends, as if the silence were a problem with him in particular, and the criticism cannot be answered on the terms it is made, because answering it would require conveying the dynamic just described, which asks for more floor and more credence than the story did. The note's word for this last part is Cassandra, and its closing sentence is that, as an autistic person, he does not have a voice in this world.

\subsection{Three levels of claim}

Three claims of different status run through the paper, and each is given its own name so that it can be assessed on its own.

The \emph{first-person account} is the note: one episode, experienced as repeated obstruction and pressure to mask, and interpreted by the person who lived it. It establishes that the experience occurred and what it was like; it does not establish how often such episodes occur or how they compare with other speakers' episodes, and nothing in the paper treats it as doing so.

The \emph{mechanism hypothesis} is the paper's immediate contribution: that in rooms where contributions are obstructed, where the account of the obstruction receives less uptake than the contributions it explains, and where the resulting silence is read as a trait of the person, masking and withdrawal are maintained by the recursion of the three, beyond what the obstruction of contributions alone would maintain. The distinctive part is the second clause. That audiences shape communicative success, that stigma encourages concealment, and that withdrawal need not mean low social interest are established findings with substantial predecessors; the claim added here is that impaired access to explanation and repair is what closes the loop, and \autoref{sec:predictions} says how that would be measured apart from the obstruction itself.

The \emph{single-root hypothesis} is the author's: that the root of autistic masking is the pervasive inability to convey the truth of one's situation, that this inability spirals into a self-reinforcing isolation, and that all autistic masking effects, and their contribution to autistic burnout, are outputs or shadows of this single abstraction. The paper keeps it as a hypothesis, states what it covers and what it does not, and gives it tests of its own (\autoref{sec:hypothesis}). The model does not establish it, and the paper does not say that it does.

Two further things are fixed before the argument begins. The vocabulary is interactional throughout: the model's quantities are rates of bids for the floor, lengths of contributions, costs of resisting, and the credence a room gives an account, and no quantity measures a person's capacity. And the hypotheses are scoped to masking and to masking's contribution to burnout. Burnout has other contributors, sensory load and life stressors among them \cite{Raymaker2020, Mantzalas2022a}, and nothing here bears on those.

\subsection{Plan}

\autoref{sec:case} reads the case with reported, inferred and hypothesized details marked. \autoref{sec:abstraction} states the structure, defines conveyance broadly enough to carry the author's claim, and places the structure against the accounts of silencing and epistemic injustice nearest to it. \autoref{sec:model} builds an illustrative model with every assumption listed and proves what it proves. \autoref{sec:masking} derives masking and its load. \autoref{sec:tower} runs the orders upward and takes apart the reply the account most often receives. \autoref{sec:hypothesis} states the two hypotheses with their boundaries. \autoref{sec:placement} places the account in the literature. \autoref{sec:predictions} specifies the studies. \autoref{sec:discussion} draws the consequences for the room. Appendix~\ref{app:note} carries the notes; Appendix~\ref{app:assumptions} tabulates the assumptions and the claims against their evidence, and records the mathematical corrections made after the first referee's brief.

\section{The case, read slowly}\label{sec:case}

\paragraph{What the note reports.} An approximate count, about six interruptions; a short window, about two minutes; a quoted redirection; a touch on the shoulder; and the author's experience of the episode as constant pressure to mask that he had to consciously resist in order to convey his thought, together with his reading of what the pressure does when a room multiplies it. That the story went unfinished is the paper's inference from the note's sequence, the redirection coming at the end of the attempt; the note does not say it outright. What follows are the paper's readings of the report, each marked for what it is.

\paragraph{The obstruction is additive, and the mechanism requires no intention.} Six interruptions in two minutes is a count, and a rate is the right description of what the count suggests, because no single interruption did the work on the note's telling; the story was lost to their accumulation. Whether any interruption was meant to silence is not inferred either way: the paper declines to infer intention, and the mechanism it proposes does not require one. This is why the note moves from the one visitor to the room full of people: several people, each interrupting at an ordinary rate, produce between them a rate at which no story survives, and the obstruction is a property of the sum. The count is not an estimate of a stable rate, and it is not necessarily six per two minutes of the author's actual floor time; it motivates the rate description and does not calibrate it.

\paragraph{The double bind.} At each interruption the speaker has two moves. He can hold the floor, which reads as inconsiderate, since the interruption was itself a bid for the floor and refusing it is a refusal of the other person; or he can yield, which loses the thought, since a story cut six times is not a story. The note's phrase is that he has to consciously resist in order to convey his thought. Both moves cost; the only move that costs nothing in the moment is not to have started. \autoref{sec:model} shows what the model makes of this: the policy that does best resists only near the end of a contribution, and when even that is not worth the attempt, silence is.

\paragraph{The floor, as the author read it.} The visitor's last utterance acknowledges the norm against interrupting and overrides it in the same breath, and it redirects him from a contribution of his own to an instrumental task on someone else's behalf. The utterance establishes nothing about her intentions; she may have thought the task urgent, or not have noticed the count. What the author read in it, and what the paper takes up as a hypothesis rather than a fact of the case, is that his turn did not have, in that room, the standing of other people's turns: that an extended contribution of his was not treated as something to be protected or, when taken, returned. A comparative claim of that kind needs observations of other speakers in the same room, which the note does not supply, and P1 in \autoref{sec:predictions} is where it would be tested. In plain terms, what the note describes is an extended contribution repeatedly disrupted and finally displaced; the conversation-analytic vocabulary of turns, transition-relevance places, overlap and floor return \cite{Sacks1974} is the right instrument for coding such episodes, and it distinguishes overlap that supports a contribution from overlap that seizes the floor, a distinction the model's bids must respect when they are measured.

\paragraph{The report is harder than the story.} The feature that makes the case more than a complaint is the last, and it is a model assumption motivated by the case rather than a fact reported in it. If asked afterward why he is quiet in company, the true answer is the dynamic above, and the dynamic is a longer turn than the rhinoceros story, told to listeners who may already hold a frame on which quietness is a trait of his. The report of the obstruction is then subject to the obstruction, at greater length and lower credence. This is what the note calls the Cassandra-like inability to convey even a fragment of the real dynamic. Cassandra's curse in the \emph{Agamemnon} is not silence; she speaks, and truly, and is not believed \cite{Aeschylus}. The case's version is exact: the one thing that would explain the silence is the thing the silence prevents from being said.

\section{The abstraction}\label{sec:abstraction}

\subsection{Conveyance, and the three conditions}

The author's claim concerns conveying the truth of one's situation, and the model of \autoref{sec:model} concerns completing a spoken contribution without losing the floor. The second is one channel of the first, and the definitions are given for the first so that the stronger hypothesis has an object.

\begin{definition}[Channel; conveyance; uptake]\label{def:conveyance}
A \emph{channel} is any route by which a person's state enters a room's record of them: an extended spoken contribution, a brief spoken request, a bodily expression such as a stim or a facial expression, a report of a need or a sensory state, a written message. A state is \emph{conveyed} on a channel when it is transmitted on it and \emph{taken up} when the room's record changes accordingly: a story understood, a request granted or at least weighed, an expression read as the state it expresses rather than as a trait of the person. A \emph{first-order} state is any state the person would have the room hold. The \emph{second-order} state attached to a failure of uptake is the account of why it failed.
\end{definition}

Uptake is not belief on demand. An account that receives fair consideration, including the listener's weighing of their own part in the difficulty, has been taken up even if the listener ends by disagreeing; an account that is matched to a prior frame and set aside has not. The distinction between completing a transmission and receiving uptake runs through the model and the predictions, since a contribution can be completed and misunderstood, and understood without any arrangement changing.

\begin{definition}[The unconveyable situation]\label{def:situation}
A person is in an \emph{unconveyable situation} with respect to a room when three conditions hold.
\begin{enumerate}[label=(\roman*)]
\item \emph{First-order obstruction.} The channels on which the person's states would be conveyed are obstructed or their contents discounted: an extended contribution is interrupted or displaced; a bodily expression is read as a trait rather than as the state it expresses; a reported need is set aside.
\item \emph{Second-order obstruction.} The account of (i) receives less uptake than the states it explains, because it asks for more floor, or more credence, or both, and because the room hears it through a frame on which the person's silence is a trait of the person.
\item \emph{Attribution.} The silence that (i) and (ii) produce is read by the room as evidence for the frame in (ii).
\end{enumerate}
\end{definition}

The three conditions are not independent symptoms; (iii) feeds (ii), and (ii) removes the route by which the person could change (i). That is what makes the structure a candidate root and not a list. The note's word for the whole is the pervasive inability to convey the truth of one's situation, where \emph{situation} does the work: what cannot be conveyed is not the story but the fact that stories cannot be conveyed. Masking, on this definition, is the suppression of a channel's content to avoid the discount in (i) or the attribution in (iii): a story shortened, a stim held still, a sensory state concealed, an expression regulated. The spoken-floor model that follows works out one channel in detail; the mapping of the others onto the three conditions is given in \autoref{sec:hypothesis}, where the stronger hypothesis is stated.

\subsection{The nearest accounts}

Three bodies of work already describe parts of this structure, and the paper's claim to add something is specific: a recursive mechanism that joins them, and tests that measure the joint.

Fricker distinguishes testimonial injustice, in which a speaker's word is given less credence than it deserves because of a prejudice about the kind of speaker she is, from hermeneutical injustice, in which a gap in the collective interpretive resources leaves a person unable to make her experience intelligible \cite{Fricker2007}. The discount in condition (ii) is testimonial where it rests on the frame. Whether condition (ii) is also hermeneutical in Fricker's sense, a gap in the shared resources rather than a particular listener's ignorance or refusal of uptake, is a further claim the paper does not make: a room that lacks the concept of room-produced silence may lack it locally, and the structural version would need the concept to be absent from the resources a culture makes available. The paper uses the testimonial reading and leaves the hermeneutical one as a question.

Dotson's account of silencing is nearer still, and the mechanism can be stated in its terms. Dotson distinguishes testimonial quieting, in which an audience fails to identify a speaker as a knower, from testimonial smothering, in which the speaker truncates her own testimony because the audience shows itself unable or unwilling to receive it competently \cite{Dotson2011}. Condition (iii) is quieting: the frame is the audience's failure to recognize the person as a knower of their own situation. The silence that the model derives in \autoref{sec:model} is smothering: the person withholds the contribution because the room's response makes attempting worse than withholding. What the present structure adds to Dotson's pair is the second order and the loop: the testimony that would explain the smothering is itself subject to quieting, and the smothering so produced feeds the quieting. Dotson's own account of the conditions under which smothering occurs, including that the audience's incompetence stems from a pernicious ignorance, should be consulted before any equivalence finer than this is claimed.

Catala, Faucher and Poirier argue that autistic people's epistemic agency is relational, constituted in part by the social environment, and introduce epistemic disablement and enablement as features of a relational setting that undermine or support it; they locate the disabling in neuronormativity and in neurotypical ignorance, particularly about autistic sociability \cite{Catala2021}. \autoref{def:situation} is a mechanism of disablement in their sense at the grain of a single room, and it is consistent with their account rather than a rival to it. The difference is in what is claimed: their account identifies the relational constitution of agency and the sources of its disablement; the present one specifies a feedback structure within that relation, derives its fixed point, and says what measurements would show the structure operating. Its novelty, so far as it has one, is the recursion and its tests, not the observation that environments silence.

\section{An illustrative model}\label{sec:model}

\subsection{The floor, the bids, and the assumptions}

An \emph{episode} is an attempt to convey an intended contribution that requires floor time $\tau > 0$. Bids for the floor on this speaker's turns arrive as a Poisson process of rate $\lambda > 0$. The \emph{exposure} is $x = \lambda\tau$, the expected number of bids over the contribution; in exposure units, bids arrive at unit rate and the contribution has length $x$. At each bid the speaker either \emph{resists}, keeping the floor at a cost $r > 0$, the social cost of reading as inconsiderate, or \emph{yields}, on which the floor passes and the contribution is lost. Attempting costs $c_0 > 0$ whether or not it succeeds, the effort of formulation and the exposure of a thought cut off; conveying the contribution is worth $V > c_0$. Silence costs nothing in the episode and conveys nothing.

The model's assumptions are listed, since what it proves is conditional on them; Appendix~\ref{app:assumptions} sorts them into simplifications and empirical hypotheses.
\begin{enumerate}[label=(M\arabic*)]
\item Bids arrive as a Poisson process, stationary within the episode.
\item With $n$ interlocutors, bids add independently: $\lambda = \sum_i \lambda_i$, and $\lambda = n\lambda_1$ in the symmetric case. This is an assumption about how a room's members behave as the room grows, not a consequence of its size.
\item A yielded bid ends the attempt; the contribution is not returned. \autoref{rem:return} relaxes this.
\item $\tau$, $V$, $r$ and $c_0$ are fixed within the episode. A contribution is treated as one stretch of floor; contributions assembled in fragments or by collaboration are outside the episode as modeled.
\item The speaker maximizes expected value within the episode, with no value assigned to what an attempt might do for future episodes. \autoref{sec:loop} says where that omission matters.
\item The menu is silence, or an attempt under any policy: a rule, possibly depending on the history of the episode, for resisting or yielding at each bid.
\end{enumerate}

The rates and costs are properties of the room's treatment of this speaker's turns, and the model's named empirical hypothesis is about them.

\begin{assumption}[Floor asymmetry]\label{prem:asymmetry}
In an allistic-majority room, an autistic speaker's extended contributions face a higher rate of floor-seizing bids, and a higher social cost of resisting them, than an allistic speaker's contributions of the same intended length.
\end{assumption}

The premise is the author's reading of the case stated as rates, and it is an empirical hypothesis: interactional research identifies the practices by which non-autistic partners block autistic contributions (\autoref{sec:placement}), and the comparative rate claim is what P1 tests. Nothing in the propositions depends on why the rates would differ.

\subsection{Completion, cost, and the policy that does best}

\begin{proposition}[Completion, cost, and the threshold policy]\label{prop:bind}
\begin{enumerate}[label=(\alph*)]
\item \emph{Stationary policies.} Under a policy that resists each bid independently with probability $\rho \in [0,1)$, the contribution completes with probability $P(\rho) = e^{-(1-\rho)x}$, the expected number of bids resisted before the attempt ends is $\frac{\rho}{1-\rho}\big(1 - e^{-(1-\rho)x}\big)$, which tends to $x$ as $\rho \to 1$, and the expected value of the attempt is
\[
U(\rho) \;=\; V e^{-(1-\rho)x} \;-\; r\,\frac{\rho}{1-\rho}\big(1 - e^{-(1-\rho)x}\big) \;-\; c_0, \qquad U(1) = V - rx - c_0 .
\]
$U$ is not in general convex in $\rho$: at $V = 10$, $r = 2$, $x = 4.6$, $U''(0) \approx -1.65$.
\item \emph{The optimal policy.} Among all policies, including those that depend on the history of the episode, the following is optimal: resist a bid exactly when the exposure remaining to completion, $y$, satisfies $y \le y^{*} := V/r - 1$; never resist if $V < r$. Its value before the cost of attempting is
\[
W(x) \;=\; \begin{cases} V - rx, & x \le y^{*},\\[2pt] r\,e^{-(x - y^{*})}, & x > y^{*},\end{cases} \qquad\text{and } W(x) = V e^{-x} \text{ if } V < r,
\]
and under it the contribution completes with probability $1$ if $x \le y^{*}$ and $e^{-(x - y^{*})}$ if $x > y^{*}$.
\item \emph{The attempt threshold.} The speaker attempts exactly when $W(x) > c_0$, that is, when $x < \bar{x}$, where
\[
\bar{x} \;=\; \begin{cases} \ln (V/c_0), & V < r,\\[2pt] (V - c_0)/r, & V \ge r \text{ and } c_0 \ge r,\\[2pt] V/r - 1 + \ln (r/c_0), & V \ge r \text{ and } c_0 < r,\end{cases}
\]
and is indifferent at $x = \bar{x}$. With $\lambda = n\lambda_1$ under (M2), this is $n < \bar{n} := \bar{x}/(\lambda_1\tau)$: the room has a size beyond which the contribution is not attempted, and the size falls as the contribution lengthens.
\end{enumerate}
\end{proposition}

\begin{proof}
(a) By thinning, the yielded bids form a Poisson process of rate $1-\rho$ in exposure units, so the attempt survives to completion with probability $e^{-(1-\rho)x}$; the resisted bids arrive at rate $\rho$ while the attempt is alive, and it is alive at exposure $t$ with probability $e^{-(1-\rho)t}$, so their expected number is $\rho\int_0^x e^{-(1-\rho)t}\,dt$. The second derivative at $\rho = 0$ is computed directly. (b) Let $W(y)$ be the optimal expected value with $y$ remaining and the floor held. Conditioning on the first bid,
\[
W(y) \;=\; e^{-y}V \;+\; \int_0^y e^{-t}\max\{0,\; W(y-t) - r\}\,dt .
\]
Suppose $V \ge r$ and substitute the stated form. For $y \le y^{*}$ every bid is resisted, since $W(y-t) - r = V - r(y-t) - r \ge V - ry - r \ge 0$, and the integral gives $e^{-y}V + (V-r)(1 - e^{-y}) - r(y - 1 + e^{-y}) = V - ry$. For $y > y^{*}$ a bid is resisted exactly when $y - t \le y^{*}$, the alternative $W(y-t) - r = r e^{-(y-t-y^{*})} - r$ being negative, and with $u = y - t$ the integral is $e^{-y}\int_0^{y^{*}} e^{u}(V - r - ru)\,du = e^{-y}(re^{y^{*}} - V)$, so $W(y) = re^{-(y-y^{*})}$; the two branches agree at $y^{*}$, where both equal $r$. The function so defined satisfies the optimality equation with the maximum attained by the stated rule, and the standard verification argument for a finite-horizon control problem with Poisson decision epochs shows that no policy, history-dependent or not, does better. If $V < r$ no bid is ever worth resisting and $W(y) = Ve^{-y}$. The completion probabilities read off the rule. (c) $W$ is continuous and strictly decreasing, so $W(x) > c_0$ holds exactly below the unique solution of $W(x) = c_0$, which is found on each branch.
\end{proof}

The proposition is the case's first two readings in the model's terms, and it corrects a statement of this paper's first draft, recorded in Appendix~\ref{app:assumptions}. The optimal policy fights for a contribution only near its end: with more than $y^{*}$ of it still to go, a bid is yielded to and the contribution lost; within $y^{*}$ of the end, every bid is resisted at a cost. A long contribution in a busy room is therefore attempted only if its first $x - y^{*}$ units can be expected to pass without any bid, which happens with probability $e^{-(x - y^{*})}$, and its last $y^{*}$ units are worth fighting through. That is the double bind as the model has it: early in the contribution, yielding loses the thought; late, resisting costs; and when the whole is worth less than the cost of beginning, silence. The two pure policies are the ends of this family, resisting throughout when $y^{*} \ge x$ and yielding throughout when $V < r$. In thirty thousand random draws of the parameters no stationary policy with $0 < \rho < 1$ beat both pure policies, but convexity fails and that observation is not a theorem; what is proved is (b). Part (c) is the giving-up, with the two features that matter later: silence is reached at a finite room size, and the size is smaller for longer contributions, so remarks survive rooms in which stories do not, and accounts of a dynamic, which are longer than stories, are the first to go. Beyond $y^{*}$ the completion probability under the optimal policy decays exponentially in exposure, and under (M2) in the number of interlocutors; the decay is in the probability of completion given an attempt, and the probability of an attempt is the separate step (c) supplies.

\begin{remark}[A returned floor]\label{rem:return}
If a yielded bid returns the floor with probability $q$, the contribution resuming where it stopped, then under a stationary policy a bid ends the attempt with probability $(1-\rho)(1-q)$, the contribution completes with probability $e^{-(1-\rho)(1-q)x}$, and the expected number of bids resisted is $\rho\,\big(1 - e^{-(1-\rho)(1-q)x}\big)/\big((1-\rho)(1-q)\big)$. Preventing bids and returning the floor are therefore different interventions: the first lowers $\lambda$, the second raises $q$, and a room can do either without the other. P4 in \autoref{sec:predictions} separates them.
\end{remark}

\subsection{The second order: the Cassandra inequality}

Let the second-order thought, the account of the obstruction, require floor time $\tau_2$, and let $\pi \in [0,1]$ be the room's degree of belief in the frame on which the person's silence is a trait of the person, a continuous summary across the room's members that the loop below updates. A contribution asks for no credence: it is \emph{received} when completed. An account asks for credence: it is received when completed and credited, and under the frame a completed account is credited with probability $\kappa \in [0,1)$, the account reading as excuse, while outside the frame a completed account is credited.

\begin{proposition}[The Cassandra inequality]\label{prop:cassandra}
In the same room and under the same stationary policy $\rho < 1$, the contribution is received with probability $P_1 = e^{-(1-\rho)\lambda\tau}$ and the account with probability
\[
P_2 \;=\; e^{-(1-\rho)\lambda\tau_2}\,\big(1 - \pi(1-\kappa)\big),
\]
so that $P_2 \le P_1$ whenever $\tau_2 \ge \tau$, with strict inequality if $\tau_2 > \tau$ or $\pi > 0$.
\end{proposition}

\begin{proof}
The account is a turn and completes by \autoref{prop:bind}(a) with $\tau_2$ in place of $\tau$; it is credited given completion with probability $(1-\pi) + \pi\kappa$. Both factors are at most the corresponding factors of $P_1$, the second being at most $1$.
\end{proof}

The inequality restates its two conditions, and that is its content: whatever makes the contribution fail, bids and length, makes the account of the failure fail at least as badly, provided the account is no shorter than what it explains and is credited no better. \autoref{sec:case} argued that both conditions hold for the account of a dynamic told to a room that holds the frame. They do not hold for every second-order move. ``Please let me finish'' is short, and in the model it is not an account at all but the act of resisting, costed at $r$ in \autoref{prop:bind}; a written message, a private conversation or a direct request to one listener are accounts on other channels, with rates, costs and credence of their own, and their availability is an empirical question the loop below treats as an exit. The useful hypothesis is therefore not that explanation is always less available than contribution but that some rooms obstruct both, and that in such rooms the account is the move the person would most need and can least make.

\subsection{The loop: persistent withdrawal in the closed model}\label{sec:loop}

Rounds $t = 1, 2, \ldots$ are episodes in one room on one channel. The room's frame weight $\pi_t$ sets the bid rate through a nondecreasing function $\lambda_t = \lambda(\pi_t)$, a room that takes the person's silence for a trait protecting the person's turns less. In round $t$ the person may attempt the contribution, worth $W(\lambda_t\tau) - c_0$ by \autoref{prop:bind}, or the account, worth $W_2(t) - c_0$ where $W_2$ is the value of \autoref{prop:bind}(b) computed at exposure $\lambda_t\tau_2$ and credited value $(1 - \pi_t(1-\kappa))V_2$, or neither; the person attempts whichever has positive value, the larger if both do. The room updates on what it receives, the correspondence bias reading an outcome as a disposition \cite{Gilbert1995}: a round in which nothing is received, whether by silence, by an attempt cut off, or by an account completed and not credited, raises the frame weight, $\pi_{t+1} = \pi_t + \alpha(1-\pi_t)$ with $\alpha \in (0,1]$; a received contribution lowers it, $\pi_{t+1} = (1-\beta)\pi_t$; and a received account lowers it by more, $\pi_{t+1} = (1-\beta_2)\pi_t$ with $\beta_2 \ge \beta$, the account being what tells the room that the silence is its own doing. The second order is thereby in the state transition: an account that completes and is discounted changes nothing, and the room's rate moves only through $\pi$.

\begin{proposition}[Persistent withdrawal]\label{prop:absorbing}
If at some round $t$ neither the contribution nor the account has positive expected value, then within the model the person attempts nothing at any later round and $\pi_t \to 1$. The state is left only by a change the model does not generate: a change of room, $\lambda(\cdot)$ falling for reasons other than the person's conveyance; a change of channel, with rates, costs and credence of its own; or a change of frame from a source other than the person's testimony.
\end{proposition}

\begin{proof}
At round $t$ nothing is received, so $\pi_{t+1} \ge \pi_t$; $\lambda$ is nondecreasing in $\pi$ and the credited value of the account is nonincreasing in $\pi$, so both expected values are no larger at $t+1$, and by induction at every later round, while $1 - \pi_{t+k} = (1-\alpha)^k(1-\pi_t) \to 0$. The exits are the quantities the proposition's hypotheses hold fixed.
\end{proof}

The proposition identifies a trap in a deliberately closed system, and it should be read as that. Within the room and channel as modeled, once neither the contribution nor its account is worth attempting, every later round confirms the frame, the frame raises the rate, and nothing the person does inside the model reverses it, because every reversal would be a conveyance and the state excludes conveyances. What the closure excludes is listed in the proposition: another room, another channel, a listener who comes to hold the situational frame from elsewhere. It also excludes, by (M5), any value the person might place on an attempt for what it does to future rounds; a person who attempts an account at negative immediate value because its success would change the room is outside the model, and whether real people do this, and whether it works, is what P5 observes. Two further things are exact. An attempt with completion probability $p$ changes the frame weight in expectation by $(1-p)\alpha(1-\pi) - p\beta\pi$, which is positive when $p$ is small: attempts likely to fail move the room toward the trap, and attempts likely to succeed move it away, so trying is not in itself an acceleration. And the exits named are mostly the room's, with one exception that is the person's: a change of channel is a move the person can make, and the model's claim is only that the move must leave the room and channel in which the loop runs.

\subsection{The fixed point: the inaccessible region, and the ghost}\label{sec:fixedpoint}

The loop has been followed for one contribution. A person's situation is many, and following the threshold of \autoref{prop:bind}(c) over all of them gives the fixed point its content: not that one story goes untold, but that a definite part of the person becomes unreachable from inside, and that the part includes its own boundary.

Let the person's \emph{situation} be a finite set $S$ of thoughts, each $s \in S$ with a floor time $\tau(s) > 0$, a value $V(s) > c_0$, and a weight $m(s) > 0$ measuring how much its conveyance matters; write $m(A) = \sum_{s \in A} m(s)$, and write $\bar{x}(v)$ for the attempt threshold of \autoref{prop:bind}(c) at value $v$, with $r$ and $c_0$ fixed. In a room of rate $\lambda$ the \emph{conveyable region} is $C(\lambda) = \{ s \in S : \lambda\tau(s) < \bar{x}(V(s)) \}$, the thoughts the person attempts there; the \emph{inaccessible region} is $I(\lambda) = S \setminus C(\lambda)$ and the \emph{inaccessible fraction} is $f(\lambda) = m(I(\lambda))/m(S)$. For $s \in I(\lambda)$ let $s'$ be its \emph{report}: the account that $s$ cannot be told in this room, and why. Two things are assumed of reports. A report is at least as long as what it reports, $\tau(s') \ge \tau(s)$, since it must carry the thought's size together with the dynamic; and a report is credited at no more than the thought's own value, $(1-\pi(1-\kappa))\,V(s') \le V(s)$, the acceptance factor of \autoref{prop:cassandra} applying to accounts, which at $\pi = 1$ reads $\kappa V(s') \le V(s)$.

\begin{proposition}[The inaccessible region]\label{prop:region}
\begin{enumerate}[label=(\alph*)]
\item $C(\lambda)$ is nonincreasing in $\lambda$ and $f(\lambda)$ nondecreasing: as the loop raises the room's rate, the conveyable region shrinks, and at the loop's stationary rate $\lambda^{*}$ the inaccessible fraction is $f(\lambda^{*})$, which is the greater the more of the situation's weight lies in long thoughts.
\item What value buys depends on the cost of resisting. For $V \ge r$ the admissible exposure $\bar{x}(V)$ grows linearly in $V$ at rate $1/r$; for $V < r$ it grows only as $\ln V$. Where a single resisted bid costs more than the thought is worth, mattering buys floor only logarithmically; where it costs less, a thought that matters more earns proportionately more floor.
\item $I(\lambda)$ is closed under reporting: if $s \in I(\lambda)$ then $s' \in I(\lambda)$. In particular the boundary of $C(\lambda)$, which thoughts can and cannot be told in this room, is itself untellable in it.
\end{enumerate}
\end{proposition}

\begin{proof}
(a) The condition $\lambda\tau(s) < \bar{x}(V(s))$ is monotone in $\lambda$ for each $s$. (b) The branches of $\bar{x}$ in \autoref{prop:bind}(c). (c) Let $s \in I(\lambda)$, so $\lambda\tau(s) \ge \bar{x}(V(s))$. The report is attempted exactly when $\lambda\tau(s') < \bar{x}$ evaluated at its credited value; $\bar{x}$ is increasing in its value and the credited value is at most $V(s)$, so that threshold is at most $\bar{x}(V(s))$, while $\lambda\tau(s') \ge \lambda\tau(s) \ge \bar{x}(V(s))$. The statement about the boundary is the case in which $s$ is any thought on it.
\end{proof}

Part (c) is the sense in which the inaccessibility is topological rather than probabilistic, and the note's word is the right one. It is not that long thoughts are unlikely to get through; a probability can be raised by trying. It is that the inaccessible region is closed under the one operation, description, that could bring any of its members within reach: the description of an unreachable thought is unreachable, and so is the description of the line between the reachable and the unreachable. A person at the fixed point therefore cannot convey the size of what they cannot convey. From outside, the region does not exist; from inside, its edge cannot be pointed at. Part (b) adds where mattering helps and where it does not: in a room where resisting is socially expensive for this speaker, which is the second half of \autoref{prem:asymmetry}, the thoughts that carry the situation's weight, the structure of a life rather than its incidents, earn floor only with the logarithm of what they are worth.

\paragraph{The ghost.} What a room holds of a person is what has been taken up, together with the frame. At the fixed point the room's record of the person is $C(\lambda^{*})$ filled out by $\pi$: the short thoughts, and the attribution that explains the absence of the rest. Call $1 - f(\lambda^{*}) = m(C(\lambda^{*}))/m(S)$ the person's \emph{presence} in the room, the fraction of the person that the room's record contains. When the inaccessible fraction is major, the person present in the room, as the room has them, is mostly not the person, and the person can see this and cannot amend it on that channel, since every amendment is a report and reports are in $I$. The relation that remains to the room is observation: the person watches the room's version of themselves as a third party would, and watches the room from a position no move of theirs on that channel can alter. This is the note's ghost, present in the room and absent from its record, viewing the outside world in the third person, and the model describes a position from which the phenomenology would be expected without adding a mechanism for it: a participant is one whose moves change the room's record, and a person none of whose moves on the channel can is, in the sense the room affords, not there. The camouflaging literature records challenges to identity and the sense of not being one's own self among masking's costs \cite{Hull2017, Cage2019}; the presence quotient locates the cost in the record and not in the person, and predicts it to track the inaccessible fraction, which is the test given as P8.

\section{Masking, its load, and burden}\label{sec:masking}

\subsection{Masking as the middle regime}

The model so far has two options, speech and silence. Masking is a third, defined here by what it does to exposure. A \emph{masked attempt} is one reshaped in advance to lower its exposure: the contribution shortened, which lowers $\tau$; delivered in the room's register with its objections pre-empted, which lowers the rate of bids on it. The forms are those the camouflaging literature catalogues under compensation, masking and assimilation \cite{Hull2017, Hull2019, Livingston2019}, and the model abstracts from the forms to their joint effect on $x = \lambda\tau$, which is one quantity whether the rate or the length is reduced. A masked attempt faces exposure $\mu x$ with $\mu \in (0,1)$ and pays an additional cost $c_M > 0$, the effort of the reshaping, on top of $c_0$; any loss of content from shortening is counted in $V$. Not every accommodation of an audience is a mask, and the model does not say otherwise; what it prices is the reshaping that the camouflaging literature's respondents describe as concealing or suppressing what they would otherwise have said. The person now chooses among silence, an unmasked attempt and a masked one, each attempt under its optimal policy.

\begin{proposition}[The forced-masking band]\label{prop:band}
Let $B_u(x) = W(x) - c_0$ and $B_M(x) = W(\mu x) - c_0 - c_M$ be the best expected values of an unmasked and a masked attempt, with $W$ as in \autoref{prop:bind}(b), and let $\bar{x}_u = \bar{x}(c_0)$ and $\bar{x}_M = \bar{x}(c_0 + c_M)$ be the exposures at which $W$ equals $c_0$ and $c_0 + c_M$.
\begin{enumerate}[label=(\alph*)]
\item The unmasked attempt has positive value exactly when $x < \bar{x}_u$, and the masked attempt exactly when $x < \bar{x}_M/\mu$. The \emph{forced band} $\bar{x}_u \le x < \bar{x}_M/\mu$ is nonempty exactly when $\mu\,\bar{x}_u < \bar{x}_M$, that is, when the mask lowers exposure by more than its cost lowers the attempt threshold; on it the masked attempt is the only speech option in the menu with positive expected net value.
\item Below $\bar{x}_u$ the person may still mask, and does so exactly when $B_M(x) > B_u(x)$.
\item The person is silent exactly when $x \ge \max\{\bar{x}_u, \bar{x}_M/\mu\}$.
\end{enumerate}
\end{proposition}

\begin{proof}
$W$ is strictly decreasing, so each attempt's value is positive on an initial interval of exposure; (a) compares the two intervals, (b) and (c) read off the values.
\end{proof}

The band is the proposition's content, and it is a statement about expected value within the menu, not about possibility: the unmasked completion probability is positive at every finite exposure, and what the band says is that in a room whose exposure lies in it, the unmasked attempt is not worth making and the masked one is, so that the person who speaks at all speaks masked. The band widens as the mask becomes more effective ($\mu$ falls) and narrows from above as it becomes more costly (a rise in $c_M$ lowers $\bar{x}_M$ and so the band's upper end, $\bar{x}_u$ being unchanged); its position in the number of interlocutors is $\bar{n}_u \le n < \bar{n}_M$ with $\bar{n}_u = \bar{x}_u/(\lambda_1\tau)$ and $\bar{n}_M = \bar{x}_M/(\mu\lambda_1\tau)$, and the three regimes are ordered in the room's size: speech, masked where the mask pays and unmasked where it does not, in small rooms; forced masking in middling ones; silence in large ones; with the boundaries moving inward as the contribution lengthens. This is the note's observation that the pressure to mask is the room's and is constant: at every exposure in the band, the choice between masking and silence is renewed at each episode, and the person who wants to be heard pays $c_M$ each time. Where resisting is cheap the band lies far out, since $\bar{x}_u$ then grows as $V/r$; the band is a feature of rooms in which resisting is expensive for this speaker, which is again the second half of \autoref{prem:asymmetry}.

\subsection{Load, burden, and burnout}

Let the person pass through episodes $t = 1, \ldots, T$ in rooms of exposures $x_t$, and let the \emph{masking load} be $L^{M}_T = \sum_{t \le T} c_M\,\mathbf{1}[\text{episode } t \text{ masked}]$. Three claims are kept apart. That masking costs effort in each episode is stipulated by the model. That the effort accumulates with repetition follows: at a constant per-episode cost the load grows linearly in the number of masked episodes, and the band's contribution to it is not a larger per-episode cost, which is $c_M$ wherever masking occurs, but the fact that in the band every contribution that is spoken at all costs it, whereas below the band only the contributions for which the mask pays do. That accumulated effort, recovery, other demands and support bear on burnout is a hypothesis with evidence of its own and not a theorem of this model.

Autistic burnout, as defined from lived experience, is a state of chronic exhaustion, loss of skills and reduced tolerance to stimulus, produced when expectations outweigh abilities under cumulative load without adequate supports \cite{Raymaker2020}, with exhaustion, withdrawal and executive difficulty as its core features \cite{Higgins2021}. In the online-post study of Mantzalas and colleagues, masking was the most commonly reported reason for burnout in that dataset \cite{Mantzalas2022b}, and their conceptual model treats it as a risk factor interacting with stressors and with barriers to support \cite{Mantzalas2022a}; their measurement study finds burnout measurement still developing and its relation to masking in need of continued investigation \cite{Mantzalas2024}. None of these validates an equation from cost to burnout, and the paper does not supply one. If a bookkeeping is wanted, the following is an unvalidated one and is identified as such: a burden $B_t \ge 0$ with
\[
B_{t+1} \;=\; \max\{0,\; B_t + c_M M_t + d_t - R_t\},
\]
where $M_t$ indicates masking in episode $t$, $d_t$ is the episode's other demand, sensory load among it, and $R_t$ is recovery; burnout, on the lived-experience definition, is the regime in which $B_t$ stays high because $c_M M_t + d_t$ exceeds $R_t$ on average. Utility costs and resource depletion need not be one quantity, and the bookkeeping is a place to put measurements, not a model of the syndrome. The note's phrase, that burnout rises exponentially at the root, is the author's; the formalization supports linear accumulation at constant cost, and a faster growth would need a feedback the model does not contain, such as burnout raising $c_M$, which the burden equation can carry as a hypothesis: a rise in $c_M$ lowers the band's upper end, so that silence arrives at a smaller room, and burnout so reached feeds the loop of \autoref{prop:absorbing}.

\paragraph{The silence load.} Masking is not the only cost the band and the region impose, and the second is the one the author's second note names: thoughts that pain the person to remain unsaid. In the model a thought $s$ not attempted forgoes its value, and the \emph{silence load} is $L^{S}_T = \sum_{t \le T} \big(V(s_t) - c_0\big)\,\mathbf{1}[s_t \in I(\lambda_t)]$, the value of what went unsaid, summed over the episodes in which it would have been said. It is zero for a thought not worth saying and largest for the thoughts that carry the situation's weight, which by \autoref{prop:region}(a) are the ones the region holds. It is a term in the burden's $d_t$ that the masking literature does not count, since nothing observable is done when it is incurred.

\section{Beyond the second order: the tower, and the conflation}\label{sec:tower}

\subsection{The tower of orders}

The second-order thought was the account of why a first-order thought failed. The construction iterates. Call $s^{(1)} = s$ and let $s^{(k+1)}$ be the account of why $s^{(k)}$ was not received: the order-three thought explains why the explanation was not credited, and so on up. Each order is an account and faces the frame, and each must carry the previous order's content together with the reason it failed, so that $\tau_{k+1} \ge \tau_k + \delta$ for some $\delta \ge 0$, the floor time a reason takes.

\begin{proposition}[The tower]\label{prop:tower}
In one room under a stationary policy $\rho < 1$, the order-$k$ thought is received, completed and credited, with probability $P_k = e^{-(1-\rho)\lambda\tau_k}\,(1-\pi(1-\kappa))^{\mathbf{1}[k \ge 2]}$, and
\begin{enumerate}[label=(\alph*)]
\item $P_{k+1} \le P_k$ for every $k$: no order is more conveyable than the one below it;
\item if $\delta > 0$ then $P_k \le P_2\, e^{-(1-\rho)\lambda\delta(k-2)}$ for $k \ge 2$, and the expected number of orders above the first that are ever received is at most $P_2/(1 - e^{-(1-\rho)\lambda\delta})$, which is finite and falls as the room grows.
\end{enumerate}
\end{proposition}

\begin{proof}
(a) The exponential factor is nonincreasing in $\tau_k$ and the credibility factor is at most $1$ and constant from the second order. (b) $\tau_k \ge \tau_2 + (k-2)\delta$ gives the bound on $P_k$, and the geometric series bounds the sum.
\end{proof}

The tower is the Cassandra inequality run upward, and it says what explanation can and cannot do on one channel in one room. A person who answers each failure of uptake with an account of that failure climbs a tower whose steps are each less reachable than the last, and the tower has a finite reach: beyond a few orders, nothing is received at all. The person cannot explain their way out on that channel, not because explanations are poor but because each explanation is a longer turn in the same room, discounted by the same frame. Other channels, and listeners who do not hold the frame, are where the tower does not apply, and \autoref{prop:absorbing} named them as exits.

\subsection{The third order: the passive example}

The note records the reply that an account of the obstruction most often receives: \emph{I have plenty of things I've never talked about with others, but you don't see it bothering me.} The reply is an event at the third order. The second-order account was heard, and its uptake failed by assimilation: the listener matched it to the one model of unsaid things they possess, their own, and on that model found nothing to be bothered by. The thought that would clear the assimilation is the order-three thought, why the listener's example is not the person's case, and \autoref{prop:tower} prices it. The note calls the reply existentially cruel, and the model says in what the cruelty consists: the reply converts the person's explanatory move into evidence for the frame. The listener has just found, in their own experience, that unsaid things do not hurt, so the person's claim that they do is read as the trait the frame already posits, and in the loop of \autoref{sec:loop} this is an account completed and not credited, which raises $\pi$. The second-order account, received this way, leaves the person worse placed than before it was made.

\paragraph{Why the conflation is a fallacy.} The listener's inference runs from their own case to the person's: unsaid things do not bother me, so unsaid things need not bother you. The inference is valid only if the two sets of unsaid things have the same structure, and the model separates them on exactly two parameters. Write $S_L$ for the listener's situation. The listener's turns are protected, so their conveyable region is their whole situation, $C_L = S_L$, and their unsaid things are $U_L = \{ s \in S_L : V_L(s) < c_0 \}$, the thoughts not worth the saying. The person's unsaid things are $I(\lambda)$, thoughts with $V(s) > c_0$ not attempted because the room's rate makes attempting worse than silence. Both sets present as things never talked about, and they differ as follows.
\begin{enumerate}[label=(\roman*)]
\item \emph{Value.} The silence cost of a thought is $(V(s) - c_0)\,\mathbf{1}[s \text{ unsaid}]$, incurred at each episode in which it would have been said. On $U_L$ it is negative or zero by construction: the listener's unsaid things are unsaid because they are not worth saying, and their silence costs nothing. On $I(\lambda)$ it is positive, and largest for the thoughts the region holds, the long ones that carry the situation. The listener's example is a case of zero silence cost offered as a model of a case of maximal silence cost.
\item \emph{The option.} The listener can convert any member of $U_L$ into a said thing by choosing to, since $C_L = S_L$; the option to speak is intact and merely unexercised. The person cannot convert any member of $I(\lambda)$ on this channel, since the attempt is not worth making at that rate: the option has been removed, in the sense the model gives the word, that no attempt has positive expected value. The listener's example is a case of an intact option offered as a model of a removed one.
\end{enumerate}
The two parameters are the note's two phrases. Things that \emph{pain} the person to remain unsaid are (i): the silence cost is the pain, and it is paid only on thoughts worth saying. The person's \emph{capacity to convey} under assault is (ii): the capacity is the conveyable region $C(\lambda)$, the room's rate shrinks it, and by \autoref{prop:region}(c) the shrinkage cannot be registered with the room. The listener's own region does not shrink and their example contains no capacity at all, only a choice.

\paragraph{Active and passive.} The note's word for the difference is that the dynamic \emph{muzzles}, and the word is exact on the model. The listener's silence is passive: it is the absence of an act, and nothing maintains it but the listener's own disinclination. The person's silence is maintained by the room, continuously: the rate $\lambda$ is the sum of the room's members' bids, an ongoing action of others, and the silence persists exactly as long as the action does. A muzzle does not change what the animal would say; it removes the saying, and it is held on from outside. The conflation takes a silence that nothing holds in place for a silence that the room holds in place, and the fallacy is the substitution of the passive for the active. This is also why the reply, offered as reassurance, lands as an obstruction: it is spoken from inside a region in which the option is intact, to a person for whom the option is gone, by someone who cannot see the option because they have never been without it.

\section{The two hypotheses, their boundaries, and what would count against them}\label{sec:hypothesis}

\begin{hypothesis}[The mechanism]\label{hyp:mechanism}
In rooms where the three conditions of \autoref{def:situation} hold, masking and withdrawal are maintained by their recursion: the discount on the account of the obstruction, and the attribution that the resulting silence feeds, sustain masking and withdrawal beyond what the obstruction of contributions alone would sustain.
\end{hypothesis}

The hypothesis is the paper's immediate contribution, and its distinctive content is the second clause: obstruction of contributions is a known pathway into withdrawal, and what is added is that impaired access to explanation and repair closes the pathway into a loop. The test is correspondingly specific, in P4: protecting contributions and giving the account fair uptake are separable interventions, and the mechanism predicts that the second has an effect over and above the first.

\begin{hypothesis}[The single root; the author's]\label{hyp:root}
Every effect classed as autistic masking, and masking's contribution to autistic burnout, is an output of the unconveyable situation of \autoref{def:situation}: the first-order obstruction, the second-order obstruction, and the attribution that joins them, acting on a person who wants to be heard.
\end{hypothesis}

The hypothesis is the note's claim that all masking effects are outputs or shadows of a single abstraction, and the model gives the word \emph{shadow} a sense: each named phenomenon is the one structure seen through one observable. The exponential decline of completion with company is the threshold policy's law. The double bind is the policy's shape. Learned silence is the attempt threshold. Masking is the band. The experience of pressure to mask is the renewal of the band's choice at each episode. Isolation is the closed loop's persistent state. The deficit attribution is the room's update. The Cassandra experience is the inequality. Masking's load is the band's per-episode cost summed. The inaccessible fraction of the situation is the fixed point's region, the untellable edge is its report-closure, the ghost is the presence quotient, the silence load is the region's forgone value, and the passive-example reply is the tower's third step. None of these needs a further cause once the structure and a person who wants to speak are given; that is what the hypothesis asserts, and the model shows that the assertion is coherent on one channel. It does not show that the assertion is true, and the model establishes neither that all withdrawal comes from the three conditions nor that every instance of masking requires all three.

\paragraph{The domains of masking.} Masking is wider than extended spoken contributions \cite{Pearson2021, Hull2019}, and the hypothesis says that each of its domains instantiates the structure on a channel of its own. The mapping is stated so that it can be tested domain by domain rather than extended after each counterexample. \emph{Suppressing a stim}: the channel is bodily expression; the first-order obstruction is a discount, the stim read as a trait rather than as the regulation it performs; the second-order account, why one stims, is heard through the same frame; the person stills the channel to avoid feeding the attribution. \emph{Concealing sensory distress}: the channel is the report of a state; the first-order obstruction is that the report is set aside or read as oversensitivity; the account of why the room's sensory conditions are different for this person is longer and less credited than the report; the person conceals the state. \emph{Regulating expression in a silent meeting}: the channel is the face; the obstruction is a discount on what the face would convey; the account would require a floor the meeting does not give. \emph{Scripting}: a mask on the spoken channel that lowers exposure at a cost, the band's instrument. In each case the testable content is that the masking should relax where the channel's content receives uptake and its account is credited, and persist where they are not, independently of other pressures. Two domains the hypothesis does not cover are named so that its boundary is visible: masking for physical safety, where stigma acts through a channel that never touches the record, and adaptations a person endorses and would keep in any room, which are not masking in the restrictive sense the camouflaging literature gives the word.

\paragraph{Acquisition and maintenance.} The model is a model of maintenance: it says what keeps masking and silence in place in a room that currently obstructs. Masking acquired under past rooms can persist in a room that no longer obstructs, as a policy learned on expectations that have not yet updated, and the mechanism predicts that it decays as the expectations do rather than at once. A person whose masking persists in a protecting room is therefore not a counterexample until the time the update would take has passed, and that time is a quantity P4 must allow for. The hypothesis is bounded, not rescued, by this: if masking persists indefinitely in rooms where both orders of conveyance are repeatedly demonstrated to succeed, the root is wider than conveyance.

\paragraph{Population and channel.} The spoken-floor model addresses people who make extended spoken contributions in rooms whose norms they do not set. Its application to children, to nonspeaking people, to users of communication aids, and to people with different support needs requires the channels of \autoref{def:conveyance} to be specified for each, and nothing here has done so.

\paragraph{What would count against them.} Against the mechanism hypothesis: fair uptake of the account, verified, with no effect on masking or withdrawal beyond that of protecting contributions. Against the single-root hypothesis: masking that persists, after the update time, in rooms where both orders of conveyance succeed on the relevant channel, which would show a root wider than conveyance; masking produced by stigma through channels that never touch the record, which bounds the hypothesis to the channels that do; and, for the model on which both rest, attempt rates falling with company while the room's rate of floor-seizing bids does not, which would refute the premise.

\section{Placement}\label{sec:placement}

\paragraph{The double empathy problem.} Milton's thesis is that the breakdown between autistic and non-autistic people is mutual, a disjuncture in reciprocity between two differently disposed social actors, and not a deficit on one side \cite{Milton2012}. The present account is downstream of it and adds the asymmetry of standing. Mutual misunderstanding explains why a story told across neurotypes may not land; it does not by itself explain why one party's contributions would be interrupted and the other's protected, nor why one party's account of the breakdown would be discounted. \autoref{def:situation} is what the double empathy problem looks like in a room where one of the two parties sets the norm of the floor, if \autoref{prem:asymmetry} holds.

\paragraph{Social motivation.} The social motivation theory of autism proposes reduced social orienting, reward and relationship maintenance as a developmental account of autistic social behavior \cite{Chevallier2012}. Jaswal and Akhtar argue that the inference from appearing to being uninterested is contradicted by autistic testimony and reinterpret the behaviors it rests on \cite{Jaswal2019}; the systematic review of camouflaging in pursuit of friendship finds autistic people camouflaging because they want friendship, withdrawing when the camouflage exhausts or fails while still wanting it, and concludes that a motivation account misreads the behavior \cite{Ridgway2025}. The model's contribution to this exchange is modest and exact. It holds the value of being heard, $V$, fixed, and derives attempts that fall with exposure while $V$ does not move; it therefore shows that a pattern of stable desire and declining attempts is consistent with an interactional account, and so challenges the inference from participation alone to desire. It does not show that desire is in fact stable, which is an assumption to be measured (P3), nor that motivation, opportunity, prior exclusion and anticipated cost do not interact; and since $\bar{x}$ grows as $V/r$ once resisting is worth it, a person who wants very much to be heard can earn a substantially larger room, so the model does not say that wanting makes no difference.

\paragraph{Interactional obstruction.} Silver and Parsons studied how non-autistic communication partners enable or block autistic people's contributions to significant conversations and found that interruptions, unsolicited ideas and unhelpful questions functioned as blocks, with the autistic person's communicative competence emerging relationally rather than residing in the individual \cite{Silver2025}. Their study identifies the practices the model's bids would be made of; it is not evidence that matched-length autistic contributions face a quantitatively higher rate of such bids than non-autistic ones, nor that rates add across interlocutors, which are the premise's content and P1's to test. Autistic adults describe group conversation as moving too fast to enter, report feeling ignored, and some withdraw into watching rather than contributing \cite{Black2023}; they describe themselves as a social minority who must conform to what neurotypical friends and family want and are used to, and find time with other autistic people easier \cite{Crompton2020b}. Non-autistic observers form less favorable first impressions of autistic adults from brief samples, on presentation rather than content \cite{Sasson2017}, which is one route by which a contribution could come to lack standing before it is heard; diagnostic disclosure, by contrast, improved first impressions in a later study \cite{SassonMorrison2019}, so the effect of a label on how an account is heard is open and is treated as such in P2. On information transfer, the first diffusion-chain study found same-neurotype chains more effective than mixed ones \cite{Crompton2020}, and its registered replication with $311$ participants found no difference between single- and mixed-neurotype chains \cite{Crompton2025}; the two are cited together, and neither measures conversational standing or interruption, which are different quantities from the transmission of a fixed message along a chain. In groups of four, all-autistic groups rated their interactions more enjoyable and friendlier than both mixed group types, and autistic adults' rapport fell when non-autistic members were present \cite{Foster2025}, which bears on the mirror prediction below and qualifies any simple symmetry.

\paragraph{Camouflaging and its costs.} The camouflaging literature describes the forms masking takes and its costs, exhaustion, threats to identity, and poorer mental health, with the reasons autistic adults give falling into conventional and relational classes \cite{Hull2017, Cage2019, Livingston2019}; the CAT-Q operationalizes compensation, masking and assimilation as separate dimensions \cite{Hull2019}. The model does not compete with the catalogue; it supplies the exposure at which reshaping is reached for and the cost it incurs per episode on one channel, and \autoref{sec:hypothesis} says how the other dimensions would map. Pearson and Rose argue that masking is a response to stigma across domains rather than a strategy freely chosen \cite{Pearson2021}, and the minority stress model finds everyday discrimination, internalized stigma and concealment predicting poorer mental health over general stress \cite{Botha2020}. The band's choice between a cost and a loss is the illusion of choice stated formally for one domain, and concealment is the band's cost named from the inside.

\paragraph{Learned helplessness.} The threshold resembles learned helplessness, in which uncontrollable outcomes produce a failure to act \cite{Maier1976}. The resemblance is partial and the difference matters. The silence of \autoref{prop:bind}(c) is a correct policy within the model, not a failure to learn control: the person has learned the room exactly, and silence is what the room's rates make optimal. What the person cannot do is convey that this is so.

\paragraph{The companion.} The companion extraction develops a cost model of interaction under governing norms, the effort an interaction extracts being a baseline plus a mismatch between a person and the conventions the interaction runs on, with a within-dyad crossover designed to test it \cite{NormCost}. The present model is the special case in which the convention is the room's norm of the floor and the mismatch is the rate at which that norm leaves the person's contributions open to bids. The companion's Principle of Symmetry, that the designations of a reference branch and another trade places under an admissible redrawing of the scheme that produced them, suggests the mirror of \autoref{prem:asymmetry} under autistic norms; it is stated below as an exploratory prediction, since a numerical majority is not a measurement of whose norms govern, and the symmetry assumption is defined there locally.

\section{The studies}\label{sec:predictions}

Each study names its question, its design, and what would weaken the claim it tests. The first three test the model's premise and its two inequalities; the fourth tests the mechanism hypothesis; the fifth and sixth test persistence and symmetry; the last two test the fixed point's consequences.

\paragraph{P1, conversational obstruction and composition.} \emph{Question:} do autistic speakers face more floor-seizing bids, and less reliable floor return, in allistic-majority groups than allistic speakers do? \emph{Design:} naturalistic recordings together with structured tasks in which the intended length and content of a contribution can be fixed before any interruption, so that intended and realized length are not confounded; group size and composition varied separately where feasible. \emph{Measures:} bids per unit of actual floor exposure, coded for floor seizure against supportive overlap; completion of the intended contribution; resumption, delay and listener understanding. \emph{Analysis:} speakers nested in groups with repeated interactions, familiarity, role, topic and status recorded; the Poisson description tested against alternatives rather than assumed, and the decay of completion with group size estimated rather than imposed. \emph{Weakens:} adequately precise evidence of no asymmetry in rates or in floor return where the premise predicts one. This would weaken the rate mechanism and leave other forms of obstruction open.

\paragraph{P2, fair uptake of self-explanation.} \emph{Question:} is a person's account of obstruction discounted beyond what its length, presentation and evidence warrant? \emph{Design:} controlled interactions in which listeners receive a speaker's contribution and, separately, the speaker's account of having been obstructed, with length varied independently, and comprehension, credence, responsibility attribution and willingness to change behavior assessed as separate outcomes, since a story and an account of an interaction are not the same kind of item. Diagnostic disclosure manipulated where relevant, with the listeners' autism knowledge measured, the label's effect being an open question in light of the disclosure study \cite{SassonMorrison2019}. \emph{Weakens:} accounts receiving fair consideration and changing the interaction in the settings where recursive obstruction is predicted. A beneficial effect of disclosure would weaken a label prediction, not the recursive hypothesis.

\paragraph{P3, desire, anticipated cost, and attempts.} \emph{Question:} can attempts to contribute fall while the desire to contribute does not, and is the fall related to obstruction? \emph{Design:} repeated within-person sampling before and after interactions, with desire, expected chance of being heard, anticipated effort, attempts made and uptake experienced measured separately; no amount of desired contact assumed, and no change in desire after repeated exclusion ruled out, since neither a preference for solitude nor a low interest in company is a deficit to be explained away. \emph{Weakens:} precise failure of expected obstruction to predict attempts once the alternatives are addressed. A joint fall in desire and attempts would not by itself identify the cause.

\paragraph{P4, change the interaction, and measure the mechanism.} \emph{Question:} do protecting contributions and giving the account fair uptake reduce unwanted masking, and does the second do so over and above the first? \emph{Design:} a counterbalanced or randomized comparison of usual interaction, explicit floor protection (returning the floor after bids, which raises $q$ in \autoref{rem:return}, and reducing bids, which lowers $\lambda$, recorded separately), explicit recognition of the speaker's account, and their combination, with an attentive control and a check that each condition changed the conversational quantity it was meant to. \emph{Measures:} completion, resumption, uptake, perceived agency, masking effort and fatigue, with the mediator measured directly; a within-episode adaptation of the CAT-Q \cite{Hull2019} treated as a new instrument requiring development, since the validated form is a trait self-report, and anticipatory dread measured on a suitable instrument rather than by description. Carryover and the delayed decay of learned masking allowed for. \emph{Weakens:} sustained, verified improvement in opportunity and uptake without the predicted reduction in unwanted masking after the update time. The interaction of the two interventions is the mechanism hypothesis's own prediction and is a hypothesis, not a result of the model. The model does not predict a larger benefit in larger rooms: under yielding the completion gain from a hazard multiplier $a$ is $e^{-ax} - e^{-x}$, which rises, peaks at $x = -\ln a/(1-a)$ and falls, so that both conditions can have low completion in highly obstructive rooms.

\paragraph{P5, persistence and recovery.} \emph{Question:} do attribution and participation form a persistent feedback process, and what changes it? \emph{Design:} longitudinal measurement in a fixed room, a family or a workplace, of participation, perceived obstruction, the room's attributions, repair attempts on every channel, masking and changes of context, with direct explanation, writing, allies, boundary-setting and changes in audience behavior observed rather than excluded. \emph{Weakens:} explanations reliably restoring participation in rooms the model counts as closed, or attributions failing to predict later obstruction. Persistence alone is not evidence for the loop, since stable exclusion, stable preferences, fatigue or fixed roles also produce stable participation; the test is whether exits are less likely under sustained obstruction and whether the room's attributions track the person's silence.

\paragraph{P6, norms, composition, and standing.} \emph{Question:} which effects follow the group's norms and composition, and which reflect unequal standing or history? \emph{Design:} compositions compared while norm-setting, authority, familiarity and partners' accommodation are measured; the mirror, that non-autistic speakers in autistic-majority rooms face the pattern of P1 and P3 in reverse, preregistered as exploratory, with the conditions under which the premise's parameters would reverse stated in advance. \emph{Weakens:} evidence against the preregistered symmetry. A failure of symmetry does not locate the cause in the autistic speaker; it identifies asymmetric environments or different mechanisms, and the rapport findings of Foster and colleagues already make simple symmetry doubtful for rapport \cite{Foster2025}.

\paragraph{P7, the conflation.} \emph{Question:} when autistic adults report obstructed conveyance, is the most frequent reply an assimilation to the listener's own voluntary reticence, and does it turn on the two parameters of \autoref{sec:tower}? \emph{Design:} vignettes varying whether the report states that the unsaid thought was one the speaker wanted to say and that the room, not the speaker, removed the option. \emph{Weakens:} a reply frequency unaffected by either statement, which would show the conflation not to turn on value and option.

\paragraph{P8, the presence quotient.} \emph{Question:} does the observer stance the note describes, watching the room's version of oneself in the third person, track the estimated inaccessible fraction within person? \emph{Design:} experience sampling with an observer-stance item alongside the measures of P3, across rooms of different size and composition and contributions of different length. \emph{Weakens:} no within-person covariation of the stance with company, with contribution length, or with floor protection.

\paragraph{The single root's own tests.} If the stronger hypothesis is to be tested and not only bounded, it needs conditions in which both contribution and fair uptake are repeatedly demonstrated on a channel while other pressures, concealment incentives among them, are varied independently, with the decay time of learned masking fixed in advance and current conditions measured apart from past ones. Successful floor interventions would support the conversational pathway; they would not establish that every masking effect has that root. A null result in one intervention would need its implementation and precision assessed before any conclusion about the root.

\section{Discussion}\label{sec:discussion}

The account relocates the maintenance of masking from the person to the room, and the relocation has a consequence for what can be done. Within the closed loop of \autoref{prop:absorbing}, the person has no move on the channel the loop runs on: every move there is a conveyance, and the state excludes conveyances. The moves the model leaves are the room's, and one of the person's. A room can protect a contribution, which lowers $\lambda$; it can return the floor after a bid, which raises $q$; it can give the account fair uptake, or hold the situational frame before the person's testimony is needed, which lowers $\pi$ and with it the discount. The person can change channel, and whether a letter, an ally or a different room reaches what the floor could not is a question the model poses and P5 observes. Agency and support need not originate in different places; what the model says is where, on one channel, the person's own agency runs out.

The note's last sentence, that as an autistic person the author does not have a voice in this world, is the author's testimony, and the model does not replace it or certify it. What the model adds is a structural reading of what the testimony describes, under the model's assumptions: a voice is a standing to be heard, a contribution that is protected and an account that is credited; and in a room at or beyond the threshold, on the channel the loop runs on, the person's contributions are not protected and the account of why is credited less than the contributions were, so that the one thing the person could say that would change the room is the thing the room is arranged not to receive. \autoref{sec:fixedpoint} gives that position its content: a region of the person closed under the one operation that could reach it, and a presence in the room that is the complement of that region. Whether a given person's rooms and channels put them there, and how far the position extends across the channels a life offers, are what the studies of \autoref{sec:predictions} would measure; the model explores the conditions under which voice becomes difficult to exercise and to recover, and the testimony reports what one such condition is like from inside.

Three limits are stated. The case is one episode, offered as the instance that fixed the structure and not as evidence for its prevalence; the prevalence is P1's and P3's to establish. The model's quantities are not measured, and its assumptions (M1) to (M6) are conveniences a measured version would relax, the single-stretch treatment of a contribution and the episode-by-episode valuation among them. And the hypotheses are scoped to masking and to masking's contribution to burnout, so that a reader who finds burnout driven in a given case by sensory load alone has found nothing against them.

\appendix
\section{The source notes}\label{app:note}
% R905 (Fable 5.1, Nov 06 2026), at the author's direction: the appendix carries the further notes as well (R905).

The author wrote the note below on November 6, 2026, at 10:30 AM, within the hour of the episode it describes, and then had Claude Fable 5.1 convert it into the present paper at his direction the same day; his further notes of the same day, from which \autoref{sec:fixedpoint} and \autoref{sec:tower} were written, are listed after it.
% R905 (Fable 5.1, Nov 06 2026), at the author's direction: the provenance statement names the further notes (R905). The case, the root claim, the double bind, the scaling by the room, the account of the silence as learned, the Cassandra formulation and the closing sentence are his, as written; the formal model, the placement in the literature and the predictions are the conversion's, written to carry his claim, and the hypothesis of \autoref{sec:hypothesis} is his claim stated in the model's terms. The note is quoted as written, with one clause corrected for grammar at the author's direction the same day; the redaction in square brackets is his own.

\begin{quote}
This is maybe the most concise root explanation of autistic masking that I’ve managed to write. (The summary is that the root itself is in the pervasive inability to convey the truth of your situation to others, which spirals into a self-reinforcing isolation—I would go so far as to argue that all autistic masking effects and their contribution to autistic burnout rising exponentially at the root are outputs or shadows of this single abstraction).  I’m noticing how [redacted person visiting the house] interrupted me about six times in a span of two minutes while I was trying to explain the funny Wittgenstein rhinoceros in the room story to my dad, and then finally at the end touched me on the shoulder and literally said “I don’t want to interrupt but why don’t you go help your sister download an app she’s looking for”. This is maybe a very good example case of why being an autistic person is a living hell: what she’s effectively doing here is keeping constant pressure on me to mask, which I have to consciously resist in order to convey my thought. Being in a room full of allistic people basically scales this dynamic up by several people, that way I just give up and learn to stay quiet. Basically it’s that allistic people make it artificially difficult to get even a word out of the Wittgenstein story or whatever it is, hence fueling a sense of isolation where me thinking of a funny story pathologically translates to the conclusion that I need to keep my mouth shut. Then I get criticized for not making friends like it’s a problem or deficit with me in particular, which again fuels the sense of isolation and Cassandra-like inability to convey even a fragment of the real dynamic. In effect, as an autistic person, I don’t have a voice in this world.
\end{quote}
% R904 (Fable 5.1, Nov 06 2026), at the author's direction: the clause "which I have to consciously resist in order to convey my thought" is the author's grammatical correction of the same day (his message: "please fix my grammar here in this one line from the entry").
% was: keeping constant pressure on me to mask, such that I have to consciously resist to convey my thought.

\paragraph{Further notes.} The author's later notes, quoted as written, each with its date.
\begin{itemize}
\item \emph{November 6, 2026, second note.} I would also very specifically argue that you reach a fixed point such that a major fraction of the truth of your situation becomes utterly inaccessible, effectively topologically so, to what you can convey, hence disassociation and the feeling that you're like a ghost who can only view the outside world in the third person. See what I mean? The reaction people give to this though is ``I have plenty of things I've never talked about with others, but you don't see it bothering me'', which is like an existentially cruel third-order obstruction because they can't see how the second order obstruction creates something more painful than their own passive example they're conflating it with! See what I mean? Like perhaps the reason the conflation is a fallacy, if I had to articulate it, is because of how this dynamic actively muzzles the autistic person, and so in the least is about things that very specifically PAIN them to remain unsaid---especially in the context of how it's their own abstract capacity to convey things itself that falls under assault by this concept. But hopefully you can push that farther based on my model than just that slight example I gave. Thanks!
\end{itemize}
% R905 (Fable 5.1, Nov 06 2026), at the author's direction: the itemized list of further notes begun with the second note of Nov 6, 2026, verbatim (its straight quotation marks set typographically).
% R906 (Fable 5.1, Nov 06 2026): Appendix A carried over unchanged from R905; Appendix B added on the first referee's brief.

\section{Assumptions, claims, and the corrections of the first revision}\label{app:assumptions}

\subsection{The assumptions, sorted}

\begin{center}\small
\begin{tabular}{p{0.3\linewidth} p{0.2\linewidth} p{0.42\linewidth}}
\toprule
Assumption & Kind & Purpose \\
\midrule
(M1) Poisson bids, stationary within an episode & simplification & gives exposure a single number and the completion law its form \\
(M2) bids add independently across interlocutors & simplification about group behavior & converts exposure into a function of room size; testable in P1 \\
(M3) a yielded bid ends the attempt & simplification & isolates loss of the thought from loss of the floor; relaxed in \autoref{rem:return} \\
(M4) fixed $\tau, V, r, c_0$; one stretch of floor & simplification & excludes fragmentary and collaborative contributions from the episode \\
(M5) episode-by-episode valuation & simplification & excludes the value of an attempt for future rounds; where it matters is stated in \autoref{sec:loop} \\
(M6) the menu: silence, or any policy & modeling choice & the propositions are about this menu; masks are added to it in \autoref{sec:masking} \\
Floor asymmetry (\autoref{prem:asymmetry}) & empirical hypothesis & the rates and costs differ by speaker in allistic-majority rooms; P1 \\
Report length and credit (\autoref{sec:fixedpoint}) & assumptions about accounts & an account is no shorter than, and credited no better than, what it explains; P2 \\
$\lambda(\pi)$ nondecreasing; the updates $\alpha, \beta, \beta_2$ & assumptions about the room & the loop; P5 \\
Masks act on exposure at a cost & modeling choice & the band; the CAT-Q dimensions map onto it in \autoref{sec:hypothesis} \\
The burden bookkeeping & unvalidated bookkeeping & a place for measurements; not a model of burnout \\
\bottomrule
\end{tabular}
\end{center}

\subsection{Claims against their evidence}

\begin{center}\small
\begin{tabular}{p{0.46\linewidth} p{0.48\linewidth}}
\toprule
Claim & Rests on \\
\midrule
About six interruptions in two minutes; the redirection; the touch; the pressure to mask as experienced & the source note \\
The story went unfinished & inference from the note's sequence \\
The author's turn lacked the standing of others' & the author's reading; a hypothesis, P1 \\
The account is longer and less credited than the story & model assumption motivated by the case; P2 \\
The threshold policy; the attempt threshold; the exponential decay of completion beyond $y^{*}$ & mathematical result on (M1)--(M6) \\
Attempt rates fall with room size & result on (M2) joined to the threshold; P1 and P3 \\
The Cassandra inequality & result on its two conditions \\
Persistent withdrawal; the exits & result in the closed loop; P5 for what happens outside it \\
The inaccessible region's report-closure & result on the two report assumptions \\
The forced band & result within the menu; a statement of expected value, not possibility \\
Linear accumulation of masking load & result at constant cost \\
Masking's relation to burnout & lived-experience definitions and first-person accounts \cite{Raymaker2020, Higgins2021, Mantzalas2022b}; conceptual model \cite{Mantzalas2022a}; measurement developing \cite{Mantzalas2024}; no equation validated \\
Obstructive practices exist & \cite{Silver2025, Black2023, Crompton2020b} \\
The rate asymmetry & hypothesis; P1 \\
The label's effect on an account's uptake & open; \cite{Sasson2017, SassonMorrison2019}; P2 \\
Same- versus mixed-neurotype transfer & \cite{Crompton2020} and its replication \cite{Crompton2025}, different findings \\
Rapport by composition & \cite{Foster2025} \\
The mechanism hypothesis & conjecture; P4 \\
The single-root hypothesis & the author's conjecture; its own tests \\
The ghost; the conflation's structure & the author's testimony and the model's reading; P7, P8 \\
\bottomrule
\end{tabular}
\end{center}

\subsection{Mathematical corrections made after the first referee's brief}

The first draft of this paper charged the cost of resisting over the whole potential interval of a contribution, $\rho\lambda\tau$ in expectation, which counts bids arriving after an earlier yield had already ended the attempt; it then proved that the expected value was convex in the resisting probability and concluded that no mixed policy beat the two pure ones. Under the stopped process the paper actually describes, the expected number of bids resisted is $\frac{\rho}{1-\rho}(1 - e^{-(1-\rho)x})$, the expected value is not convex (at $V = 10$, $r = 2$, $x = 4.6$, its second derivative at $\rho = 0$ is about $-1.65$), and a policy that depends on the history of the episode can beat both pure policies: with $V = 10$, $c_0 = 1$, $r = 2$ and $x = 4.6$, always yielding is worth about $-0.90$ and always resisting $-0.20$, while yielding during the first $0.6$ units of exposure and resisting through the remaining $4$ is worth about $+0.10$. The brief's correction and example are both right. \autoref{prop:bind} now gives the stationary-policy accounting with the corrected cost, proves by dynamic programming that the optimal policy among all policies is to resist exactly when the remaining exposure is at most $V/r - 1$, with the example as its instance, and derives the attempt threshold from that policy's value. The completion law under a stationary policy was unchanged by the correction; the threshold $\theta^{*}$ of the first draft is replaced by $\bar{x}$, which coincides with it only when $V < r$ or $c_0 \ge r$. Two further statements were withdrawn: that masking's contribution to burnout is largest in the forced band, since the per-episode cost of masking is the same wherever it occurs, and that a floor intervention's benefit is larger in larger rooms, since the completion gain from a hazard multiplier rises and then falls in exposure. The second-order process is now in the loop's state transition, an account completed and not credited leaving the room unchanged or confirming the frame; the exits from persistent withdrawal include a change of channel, which is the person's move; and the statement that attempting accelerates the trap was replaced by the sign condition on the expected change of the frame weight. The verification of the threshold policy by discretized dynamic programming and by simulation accompanies the paper's checks.

\printbibliography

\end{document}
